\documentclass[12pt]{article}
\usepackage[utf8]{inputenc}
\usepackage[T1]{fontenc}
\usepackage{hyperref}
\title{Varisi Language Resources for K-12 Curriculum Development}
\author{Compiled by Claude Sonnet 5 (claude-sonnet-5) for Andr\'as Kornai}
\date{}
\begin{document}
\maketitle
\section{Introduction}
Varisi (ISO 639-3 vrs; Glottocode vari1239; also Varese, Ghone) is an Oceanic language of Choiseul Province, Solomon Islands, with about 5,200 speakers (1999 figure). Very little was found: scripture and audio, but no dictionary, grammar or school materials.
\section{Primary Educational Resources}
None found in the search results.
\section{Linguistic and Reference Materials}
\begin{itemize}
\item Wikipedia overview: \url{https://en.wikipedia.org/wiki/Varisi_language}
\item Verbix page with a short verb list: \url{https://docs.verbix.com/Languages/Varisi}
\end{itemize}
No dictionary or grammar was found.
\section{Religious and Community Resources}
Scripture Earth lists Varisi scripture, including a New Testament: \url{https://www.scriptureearth.org/00eng.php?iso=vrs}. Joshua Project people-group page: \url{https://joshuaproject.net/people_groups/15775}
\section{Digital and Audio Resources}
Global Recordings Network lists audio recordings: \url{https://globalrecordings.net/en/language/vrs}
\section{Curriculum Development Context}
Varisi is used as a first language by the whole community, but no school use was found. The New Testament and the Global Recordings audio are the only substantial native-language materials located, and any curriculum would need new community-led materials.
\end{document}